\section{A signed counting tensor}\label{sec:model}

Fix a characteristic $\chi$ admitting ports, in the sense of
Section~\ref{sec:prices}. The bounded-degree expansion property can
equivalently be tested at power-of-two orders
\cite[Definition~\ext{G}{def:ports} and the discussion there]{oaG}.
In this case the port identities compress a tensor network while
preserving its contraction. We encode the signed counts of a universal
machine as coefficients of a polynomial. For exact networks the
contradiction comes from making a machine recover its own count and
add one. Extending the argument to the closure of circuit outputs gives
a stretched-exponential circuit lower bound for Hamiltonian cycle and,
when $\chi\ne2$, for the permanent.

Throughout Part~\ref{part:port}, logarithms are to base two. A wire
with $2^b$ possible labels has width $b$ bits. Unless specified to be
absolute, exponents in the length estimates may depend on $\chi$ and
the port constants.

\subsection{The coefficient of a computation}

Fix a self-delimiting binary encoding of multitape Turing-machine
programs, for example by transition tables with explicit tapes and
a fixed alphabet. An input $x\in\{0,1\}^n$ begins with the code of a
program. Run the program with read access to all of $x$. At each
step it may request a fresh private bit and branch on its value.
Give each branch $n^{100}$ simulated steps. A valid halt returns
$+1$, $-1$ or $0$, and a malformed program, an error or a timeout
contributes zero. The encoding admits deterministic universal
simulation with polynomial overhead.

Let $U_n^*(x)$ be the sum of these outputs in $\mathbb Z$. This is a
clocked universal GapP function \cite{GapP}. Viewing its table as a
tensor with $n$ bit indices, form the polynomial
\begin{equation}\label{eq:signed-witness}
  P_n^*(Z)=\sum_{x\in\{0,1\}^n}U_n^*(x)\prod_{i=1}^n Z_{i,x_i}.
\end{equation}
The choice of one variable from each pair $(Z_{i,0},Z_{i,1})$
records the bits of $x$. Thus the coefficient of the corresponding
monomial is the signed count for that input, and every monomial
has degree $n$. We use the same notation over a field after mapping
the integer coefficients to it.

To see uniform $\VNP$ membership, supply $x$ and a string of $n^{100}$
guess bits as Boolean summation variables. A deterministic verifier
feeds the bits to the program as it requests them and requires every
unused bit to be zero. Each computation path therefore has exactly
one padded guess string, and $|U_n^*(x)|\le 2^{n^{100}}$. Implement
the verifier by a uniform polynomial-size Boolean circuit of bounded
fan-in. Introduce a sum bit for each gate output and multiply the
truth-table predicates of all gates. Each predicate is a
constant-degree integral polynomial, obtained by summing products of
literals for the allowed rows. For fixed input and guesses there is
one consistent gate assignment. Weight it by the positive-output
indicator minus the negative-output indicator and by
\[
                 \prod_i\big((1-x_i)Z_{i,0}+x_i Z_{i,1}\big).
\]
The resulting polynomial has polynomial size and degree before the
Boolean sum. The same construction over each field yields
\eqref{eq:signed-witness}.

\subsection{Batches and border circuits}

A bit network contracts local tensors by summing products over paired
bit indices. A ternary equality tensor can pass the same bit to
several factors. The \emph{volume} counts tensors and wires, including
one wire for each unpaired exterior index. A network is
\emph{trivalent} if each tensor has at most three incident indices,
including exterior ones. Loops, parallel wires and disconnected
networks are allowed.

A batch of $k$ copies has the exterior table
\[
        (x^{(1)},\ldots,x^{(k)})\longmapsto
              \prod_{j=1}^k U_n^*(x^{(j)}).
\]
A network realizing this table may connect indices from different
copies internally. The first lower bound controls the cost per copy
even with this freedom.

\begin{theorem}[Exact batches]\label{thm:batches}
Fix a finite field $F$ of characteristic $\chi$.
There are constants $\sigma,\delta>0$ such that, for all sufficiently
large $n$ and every $1\le k\le 2^{n^\sigma}$, every trivalent bit
network over $F$ realizing the batch $(U_n^*)^{\otimes k}$ has
volume greater than $kn^{1+\delta}$.
\end{theorem}

\begin{theorem}[Border lower bound in port characteristics]\label{thm:stretch}
There are $\sigma,\varepsilon>0$ such that over every field $F$ of
characteristic $\chi$, including the case $\chi=0$,
\[
   \bsz_F^{[n]}(P_n^*)>2^{n^\sigma}
       \qquad\text{for all sufficiently large }n .
\]
For all sufficiently large $N$ in this characteristic,
\[
   \bsz_F^{[N]}(\HC_N)>2^{N^\varepsilon}.
\]
The same bound holds for $\Per_N$ when $\chi\ne2$.
\end{theorem}

\paragraph{How the computation uses its own count.}
Pin the program bits of $R$ tensors $U_n^*$ and join pairs of their
other bit indices. Each tensor can then check a portion of a
computation or carry out a private job. Summing the shared bits
multiplies the local path counts and adds over consistent public
computations. For an appropriate wiring, the resulting scalar counts
the successes of a prescribed program whose code and advice are
supplied as pins.

Suppose the batch in Theorem~\ref{thm:batches} had a small exact
network. Use it in place of each group of $k$ processors. An evaluator
can read the program and advice as data, recover the pins and links,
and count assignments in the substituted network modulo the
characteristic. This evaluator needs too many public steps for the
original processors to run directly. The port identities reduce the
number of bits that must be guessed and checked jointly. The remaining
checks split into short private jobs that individual processors can
execute. The original program can consequently run the compressed
evaluator on its own description and contribute one additional
success. The two evaluations of its count then differ by one in the
prime field. Section~\ref{sec:batches} gives the exact argument, whose
batch lower bound also enters the analysis of matrix multiplication
in \cite{oaM}.

For Theorem~\ref{thm:stretch}, the proposed small circuit supplies a
table of coefficient queries in formal scalar parameters. Many such
queries can again be counted together. The target coefficient vector
may, however, lie only in the closure of the circuit outputs. On a
curve reaching that vector the parameters can have poles, even
when the contracted count has a well-defined value. The Cartier
operator on differentials yields matrices for the monomials in the
parameters. Their weighted sum acts on one common nonzero vector by
the boundary value of the count. This identity lets the program add
one at a nonzero coordinate without having to compute the vector.
The projections to Hamiltonian cycle and the permanent then complete
Theorem~\ref{thm:main}.
